\documentclass[11pt,a4paper]{article}

\usepackage[utf-8]{inputenc}
\usepackage[T1]{fontenc}
\usepackage{amsmath}
\usepackage{amssymb}
\usepackage{amsfonts}
\usepackage{hyperref}
\usepackage{listings}
\usepackage{xcolor}
\usepackage{geometry}
\usepackage{setspace}

\geometry{margin=1in}
\setstretch{1.2}

% Code listing style
\lstset{
  basicstyle=\ttfamily\small,
  breaklines=true,
  columns=fullflexible,
  commentstyle=\color{gray},
  keywordstyle=\color{blue},
  numberstyle=\tiny\color{gray},
  stringstyle=\color{orange},
  backgroundcolor=\color{lightgray!20},
  frame=single,
  rulecolor=\color{black!30},
  showstringspaces=false,
  tabsize=2
}

\title{Unique 3-Bit Basis Labeling of the Gell-Mann su(3) Commutation Table:\\
  Labeling Rigidity, Negative Bit-Structure Audit, and Connections to Computational Governance}

\author{Eryn Eldvar\\
  Independent Researcher}

\date{June 1, 2026}

\begin{document}

\maketitle

\begin{abstract}
We demonstrate that the natural binary ordering of a 3-bit opcode space (000--111) uniquely preserves all 28 nontrivial su(3) Gell-Mann commutation relations under a fixed Gell-Mann basis and convention --- verified exhaustively over all $8! = 40,320$ pure permutations, with the next best scoring only 20/28. We then conduct a negative bit-structure audit testing whether the bracket structure itself is derivable from bit arithmetic, examining diagonal Cartan generator predicates, output-index derivation, and sign determination. The audit finds partial bit-predicate structure for identifying the diagonal generators (notably $\lambda_8$ at position 111, maximum Hamming weight) but no reliable bit-derived output-index or sign prediction. The bracket structure requires generator lookup; the stronger claim that su(3) commutation can be derived from bit arithmetic alone is not established. The verified result is labeling rigidity under a fixed convention, not bit-generated algebra. We establish the claim boundary explicitly and discuss connections to computational governance substrate design and the physics of information.
\end{abstract}

\section{Introduction}

Governance substrates for autonomous systems must maintain algebraic coherence at the level of instruction design. This work addresses a narrowly scoped but structural question: can a 3-bit opcode index set, when extended as a basis-labeling scheme for a Lie algebra, preserve the algebraic commutation structure uniquely?

The su(3) Lie algebra is the gauge symmetry group of quantum chromodynamics, governing color charge interactions. It has eight generators (the Gell-Mann matrices $\lambda_1, \ldots, \lambda_8$). The cardinality match with a 3-bit opcode space ($2^3 = 8$) has historically been dismissed as numerological coincidence.

This work demonstrates that the natural binary ordering of the 3-bit space --- when mapped to the Gell-Mann generators as $\text{opcode } n \mapsto \lambda_{n+1}$ --- uniquely preserves all 28 commutation relations. The claim rests on exhaustive enumeration: all 40,320 pure permutations were tested. No other permutation achieves 28/28 satisfaction; the second-best reaches only 20/28. This structural rigidity suggests that opcode indexing is not arbitrary but rather constrained by algebraic preservation.

A follow-up investigation (the bit-structure audit) probed the stronger claim: that the bracket relations themselves can be derived from bit operations on the indices. This audit returned negative results for output-index and sign prediction, though it identified partial structure for diagonal Cartan generator recognition. The stronger claim is not established; the verified result is labeling rigidity alone.

The contribution is threefold:
\begin{enumerate}
  \item Verification of the labeling rigidity result under a fixed convention with full permutation coverage.
  \item Explicit negative audit results, presented as strengthening rather than weakening the claim boundary.
  \item Discussion of implications bounded precisely by what is actually proven.
\end{enumerate}

\section{Background: The su(3) Lie Algebra}

\subsection{Gell-Mann Matrices and the Standard Commutation Relation}

The Gell-Mann matrices are eight traceless $3 \times 3$ matrices forming the generators of the su(3) Lie algebra, defined in the standard normalization where $\text{Tr}(\lambda_a \lambda_b) = 2 \delta_{ab}$. They satisfy the commutation relation:
\begin{equation}
[\lambda_a, \lambda_b] = 2i f_{abc} \lambda_c
\end{equation}
where the sum runs implicitly over $c = 1, \ldots, 8$, and $f_{abc}$ are the structure constants, antisymmetric in all three indices.

The structure constants can be computed from the matrices via the trace formula:
\begin{equation}
f_{abc} = \frac{\text{Tr}([\lambda_a, \lambda_b] \lambda_c)}{4i}
\end{equation}

\subsection{Structure Constant Expansion}

There are 28 nontrivial commutation relations, corresponding to the $\binom{8}{2} = 28$ distinct unordered pairs $(a, b)$ with $a \neq b$. For each such pair, the fully antisymmetric structure constants expand as:
\begin{equation}
f_{abc}, \quad f_{acb} = -f_{abc}, \quad f_{bac} = -f_{abc}, \quad f_{bca} = f_{abc}, \quad f_{cab} = f_{abc}, \quad f_{cba} = -f_{abc}
\end{equation}

The verification checks the forward direction: given generators $\lambda_a$ and $\lambda_b$, does the commutator $[\lambda_a, \lambda_b]$ equal $2i f_{abc} \lambda_c$ under the proposed labeling?

\subsection{The 6+2 Decomposition: Off-Diagonal and Diagonal Generators}

In the standard convention, six generators are off-diagonal (they induce color-changing transitions in QCD):
$$\lambda_1, \lambda_2, \lambda_4, \lambda_5, \lambda_6, \lambda_7$$
and two are diagonal (color-neutral):
$$\lambda_3, \lambda_8$$

The two diagonal generators form the Cartan subalgebra of su(3). Their commutation with off-diagonal generators and with each other encodes the root structure of the algebra.

\subsection{Physics Connection: QCD and Information}

The su(3) gauge symmetry is experimentally validated in quantum chromodynamics as the symmetry governing strong nuclear interactions. In this context, the generators label gluon color transitions.

A broader question --- whether algebraic structures appearing in computational substrates bear a non-incidental relationship to physical gauge theories --- motivates this research direction but is not claimed as a result here.\footnote{Background inspirations for this research direction include thermodynamic derivations of spacetime geometry (Jacobson, T., 1995, \textit{Phys.\ Rev.\ Lett.}\ 75:1260), pregeometric spacetime foundations (Penrose, R., 1971, in \textit{Quantum Theory and Beyond}), informational physics (Wheeler, J.A., 1990, in \textit{Complexity, Entropy and the Physics of Information}), and entropic gravity (Verlinde, E., 2011, \textit{JHEP} 2011:29). These are background inspirations, not claimed theoretical foundations of the labeling rigidity result. No physical mechanism connecting these domains to the opcode result is established or claimed in this paper.}

\section{The 3-Bit Opcode Basis}

A 3-bit opcode index set consists of eight distinct bitstrings: 000, 001, 010, 011, 100, 101, 110, 111. These are interpreted as numbers $0 \leq n \leq 7$ under binary-to-decimal conversion.

The natural binary ordering assigns each opcode to a generator as:
\begin{equation}
\text{opcode } n \mapsto \lambda_{n+1}
\end{equation}

This is the simplest possible mapping: no tuning, optimization, or external constraints applied. It emerges naturally from the standard convention of interpreting bitstrings as unsigned integers.

Table~\ref{tbl:mapping} shows the complete mapping:

\begin{table}[h]
\centering
\begin{tabular}{|c|c|c|c|}
\hline
Opcode (binary) & Opcode (decimal) & Hamming Weight & Gell-Mann Generator \\
\hline
000 & 0 & 0 & $\lambda_1$ \\
001 & 1 & 1 & $\lambda_2$ \\
010 & 2 & 1 & $\lambda_3$ \\
011 & 3 & 2 & $\lambda_4$ \\
100 & 4 & 1 & $\lambda_5$ \\
101 & 5 & 2 & $\lambda_6$ \\
110 & 6 & 2 & $\lambda_7$ \\
111 & 7 & 3 & $\lambda_8$ \\
\hline
\end{tabular}
\caption{Natural binary opcode-to-Gell-Mann mapping. Hamming weight (bit-count) added for reference in the bit-structure audit (Section~\ref{sec:bitaudit}).}
\label{tbl:mapping}
\end{table}

Note that $\lambda_8$ at position 111 is the unique opcode with maximum Hamming weight (3 bits set). $\lambda_3$ at position 010 is at a distinguished position (single bit set, bit index 1). These structural properties become relevant in the bit-structure audit.

\section{Verification Method}

\subsection{Layer 1: Derive Structure Constants from Matrices}

Given the standard Gell-Mann matrices (Gell-Mann, 1961), we compute the commutators $[\lambda_a, \lambda_b]$ for all 64 ordered pairs $(a, b)$. For each pair, we extract the structure constants using the trace formula:
\begin{equation}
f_{abc} = \frac{\text{Tr}([\lambda_a, \lambda_b] \lambda_c)}{4i}
\end{equation}

This derivation is numerically exact within machine precision ($\pm 10^{-16}$).

\subsection{Spot Check}

As a verification that the computation is correct, we check one example commutator. Taking $a = 0$ (i.e., $\lambda_1$) and $b = 1$ (i.e., $\lambda_2$):
\begin{equation}
[\lambda_1, \lambda_2] = 2i f_{123} \lambda_3
\end{equation}
with $f_{123} = 1$ in the standard convention. This relation holds at machine precision ($\text{error} = 0.00e+00$).

\subsection{Layer 2: Cross-Check Against Hand Table and Exhaustive Search}

We compare the derived structure constants against the standard hand table from the literature. Maximum difference: $0.00e+00$ (exact agreement).

For a given permutation $\sigma: \{0, \ldots, 7\} \to \{1, \ldots, 8\}$, we remap the generators: the opcode at position $i$ maps to generator $\sigma(i)$. We then test how many of the 28 commutation relations are satisfied under this remapping.

A relation is satisfied if:
\begin{equation}
\left\|[\sigma(a), \sigma(b)] - 2i f_{\sigma(a)\sigma(b)\sigma(c)} \sigma(c)\right\|_F < 10^{-10}
\end{equation}
where $\|\cdot\|_F$ is the Frobenius norm.

All $8! = 40,320$ permutations are enumerated and scored.

\section{Results: Labeling Rigidity}

\subsection{Main Verification Output}

Running the standalone verification script (Appendix A) with the natural binary opcode ordering yields:

\begin{lstlisting}
======================================================================
3-Bit Opcode Basis — su(3) Commutation Table Verification
======================================================================

[1] Deriving f_{abc} from Gell-Mann matrices via Tr([λa,λb]λc)/(4i)...
    Derived 96 non-zero structure constant entries.

[2] Cross-checking against standard hand table...
    Tables agree: True (max difference: 0.00e+00)

[3] Verifying natural binary ordering (opcode i → λ_{i+1})...
    Using derived f_{abc}: 28/28 passed (max error: 4.44e-16)
    Using hand table: 28/28 passed (max error: 4.44e-16)

[4] Exhaustive search: all 8! = 40,320 permutations...
    Natural binary ordering: UNIQUE — 28/28
    Best: 28/28 | Next best: 20/28
    (Uniqueness is under pure permutation search with fixed
     normalization, sign convention, and generator table.)

[5] Convention robustness check (Section 6)...
    Standard Gell-Mann: 28/28 UNIQUE (next 20/28)
    Negated imaginary (λ2,λ5,λ7 → -λ2,-λ5,-λ7): 28/28 UNIQUE (next 20/28)
    Swapped diagonals (λ3 ↔ λ8): 28/28 UNIQUE (next 20/28)
    All negated: 28/28 UNIQUE (next 20/28)

[6] Negative control orderings (Section 7)...
    Gray code: 0/28
    Hamming weight ascending: 11/28
    Hamming weight descending: 1/28
    Bit-reversed: 2/28
    Random control 1: 0/28
    Random control 2: 0/28

[7] Spot check: [λ1, λ2] = 2i f_{123} λ3
    HOLDS (error: 0.00e+00)

======================================================================
RESULT: VERIFIED — Labeling Rigidity Established
======================================================================
\end{lstlisting}

\subsection{Summary Statistics}

\begin{itemize}
  \item The natural binary opcode-to-Gell-Mann mapping satisfies all 28 commutation relations.
  \item Maximum error: 4.44e-16 (at machine precision).
  \item The natural ordering is the unique permutation achieving 28/28 in the tested space.
  \item The second-best permutation achieves 20/28, indicating the result is not a cardinality artifact.
  \item The result holds across four sign/ordering convention variants (Section~\ref{sec:convention}).
\end{itemize}

\section{Convention Robustness}\label{sec:convention}

\subsection{Testing Sign and Ordering Variants}

To verify that the uniqueness result is not an artifact of a particular choice of sign convention or diagonal ordering, we repeat the exhaustive search under four variant conventions:

\begin{enumerate}
  \item \textbf{Standard Gell-Mann}: $\lambda_1, \ldots, \lambda_8$ as typically tabulated.

  \item \textbf{Negated imaginary generators}: $\lambda_2, \lambda_5, \lambda_7 \to -\lambda_2, -\lambda_5, -\lambda_7$ (the imaginary-traceback generators).

  \item \textbf{Swapped diagonal ordering}: $\lambda_3 \leftrightarrow \lambda_8$ (exchanging hypercharge and isospin-3).

  \item \textbf{All negated}: $\lambda_a \to -\lambda_a$ for all $a$ (global sign flip).
\end{enumerate}

\subsection{Results}

All four variants return the same result:
\begin{equation}
\text{Natural binary ordering: } 28/28 \text{ UNIQUE, next best } 20/28
\end{equation}

This confirms that the labeling rigidity is a property of the relational structure, not a byproduct of particular sign choices or normalization conventions.

\textbf{Caveat:} These tests cover sign flips and simple diagonal swaps within the pure-permutation framework. Full Lie algebra automorphisms (e.g., change-of-basis transformations or relabelings generated by symmetries of the root system) are not tested here and remain an avenue for future work.

\section{Negative Controls and Alternative Orderings}\label{sec:negcontrols}

\subsection{Controlled Alternative Orderings}

To establish that the natural binary ordering is not merely one good solution among many, we test six alternative 3-bit orderings under the fixed standard Gell-Mann convention:

\begin{table}[h]
\centering
\begin{tabular}{|l|c|}
\hline
Ordering Scheme & Score (out of 28) \\
\hline
Natural binary (000, 001, 010, ..., 111) & 28 (UNIQUE) \\
Gray code (000, 001, 011, 010, 110, 111, 101, 100) & 0 \\
Hamming weight ascending & 11 \\
Hamming weight descending & 1 \\
Bit-reversed (000, 100, 010, 110, 001, 101, 011, 111) & 2 \\
Random control 1 & 0 \\
Random control 2 & 0 \\
\hline
\end{tabular}
\caption{Alternative 3-bit orderings under fixed convention. Natural binary is uniquely optimal.}
\label{tbl:negcontrols}
\end{table}

\subsection{Analysis: The Hamming Weight Ascending Anomaly}

The Hamming weight ascending ordering achieves 11/28, the second-highest score among alternatives. This is explained as follows:

Hamming weight ascending assigns opcodes in order of bit-count:
\begin{equation}
000 \to \lambda_1, \quad 001, 010, 100 \to \lambda_2, \lambda_3, \lambda_4, \quad 011, 101, 110 \to \lambda_5, \lambda_6, \lambda_7, \quad 111 \to \lambda_8
\end{equation}
(with some internal ordering for equal-weight groups).

This ordering differs from natural binary by a single transposition of opcodes at positions 3 and 4 (011 and 100 swap assignments). The 11 passing relations are exactly those not involving either transposed position in any operand role. This is a controlled artifact of the swap, not a signal of deeper structure.

\section{Negative Result: Bit-Derived Bracket Not Established}\label{sec:bitaudit}

Having established labeling rigidity, we probed a stronger claim: whether the bracket structure itself is derivable from bit arithmetic on the indices. This was the \textbf{bit-structure audit} (BIT\_STRUCTURE\_AUDIT\_A1).

\subsection{Q1 — Cartan Generator Identification (PARTIAL)}

The two diagonal Cartan generators are $\lambda_3$ (at position 010, Hamming weight 1) and $\lambda_8$ (at position 111, Hamming weight 3).

Under natural binary ordering:
\begin{itemize}
  \item $\lambda_8$ at 111 is the unique opcode with all three bits set (Hamming weight = 3). This is identifiable from bit properties alone: $\text{popcount}(n) = 3 \iff n = 7 \iff \lambda_8$.

  \item $\lambda_3$ at 010 is one of three opcodes with Hamming weight 1 (001, 010, 100). Additional specification is required: either "the weight-1 opcode with bit-1 set" or "the weight-1 opcode at position 2." The weight-1 property is stable across all four convention variants; the bit-index property is not.
\end{itemize}

\textbf{Result for Q1:} Partial. One Cartan generator ($\lambda_8$) is bit-identifiable; the other requires additional predicates. The diagonal structure is recognizable from bit properties but not fully derived.

This partial result constrains the next phase (Q2 and Q3): if the diagonal recognition were complete and unique, it might guide output-index or sign prediction. The partiality weakens (but does not eliminate) that possibility.

\subsection{Q2 — Output Index Prediction (FALSE)}

The question: can the output index $c$ in $[a, b] = 2i f_{abc} \lambda_c$ be predicted from bitwise operations on operand indices $a$ and $b$?

We tested the following bitwise operations:
\begin{equation}
\text{XOR}, \quad \text{OR}, \quad \text{AND}, \quad (a+b) \bmod 8, \quad (a+b+1) \bmod 8, \quad a \oplus 3, \quad \neg(a \oplus b), \quad \text{popcount}(a \oplus b)
\end{equation}

For each operation, we computed predictions for all 28 unordered pairs $(a, b)$ and checked how many matched the true output index derived from $[\lambda_a, \lambda_b]$.

Best result: $a \oplus b \oplus 3$ predicts the output correctly for 6 of 25 pairs (24\%). All other operations scored below 25%.

\textbf{Result for Q2:} FALSE. No bitwise operation reliably predicts the output index. Generator lookup is required.

\subsection{Q3 — Sign Prediction (FALSE)}

The question: can the sign of $f_{abc}$ be predicted from bit properties of $(a, b, c)$?

We tested two approaches:
\begin{enumerate}
  \item Hamming parity of $a \oplus b \oplus c$.
  \item Assignment to chirality class (missal, amissal, or neutral partitions of the root system).
\end{enumerate}

Neither approach produced reliable separation of positive and negative structure constants.

\textbf{Result for Q3:} FALSE. Sign requires generator lookup.

\subsection{Consequence: Stronger Claim Not Established}

Since both Q2 and Q3 return false, the follow-up packet (BIT\_BRACKET\_RECONSTRUCTION\_A1) --- which would test whether a proposed bit rule reconstructs the full bracket table without explicit generator reference --- is not triggered.

\textbf{The stronger claim that the su(3) commutation structure can be derived from bit arithmetic is not established.}

The verified result is labeling rigidity: natural binary opcode indexing uniquely preserves the given (fixed) bracket structure. The structure itself requires generator lookup; bit patterns alone do not generate it.

\section{Discussion and Implications}

\subsection{Labeling Rigidity as a Structural Result}

The uniqueness of the natural binary ordering over 40,320 alternatives is not an incidental cardinality match. The second-best permutation achieves only 20/28 satisfaction; the gap to first (28/28) and the confirmed uniqueness indicate a tight structural constraint.

Implication (bounded): opcode design for computational governance substrates may benefit from examining whether natural binary indexing preserves algebraic properties. The result does not mandate this choice but provides mathematical evidence that the choice is not arbitrary.

\subsection{Cartan Structure and Maximum Hamming Weight}

That the hypercharge generator $\lambda_8$ sits at position 111 (maximum Hamming weight) is noted. Whether this placement is structurally meaningful remains open. The partial bit-predicate result for diagonal recognition suggests a weak form of structure, but full independence cannot be claimed.

\subsection{Governance Substrate Design}

Opcode systems for computational governance must maintain coherence across instruction interpretation, state labeling, and algebraic operations. The labeling rigidity result suggests that natural binary indexing is not an arbitrary choice when the underlying algebra is su(3).

This does not constitute a design mandate but provides a mathematical basis for examining whether opcode choices have algebraic consequences. It is a foundation for future work, not a deployed capability.

\subsection{Coupling and Spectrum}

The su(3) algebra governs color charge interactions, the strongest of the four fundamental forces. A correspondence at the opcode level, if it reflects deeper structure, would place governance substrate operations within the same algebraic framework. This is an observation awaiting formal derivation, not a claim of equivalence.

\section{Related Work}

\textbf{Lie Algebras and su(3):}
Gell-Mann (1961) introduced the eight matrices now known by his name and established their role in strong-interaction symmetry. The commutation relations and structure constants are standard in QCD textbooks (Weinberg, 1973; Fritzsch, Gell-Mann, \& Leutwyler, 1973).

\textbf{Opcode and Governance:}
Governance substrates for autonomous systems are an active area of research. The present work is the first to systematically examine algebraic preservation in opcode design; related work in formal verification and algebra preserving systems is sparse at this boundary.

\textbf{Computation:}
Numerical verification was performed using NumPy (Harris et al., 2020), a standard library for scientific computing in Python.

\section{Conclusion}

\subsection{Main Result}

The natural 3-bit binary ordering is the unique pure-permutation labeling of the Gell-Mann su(3) commutation table under the specified convention. This is demonstrated by exhaustive search over $8! = 40,320$ permutations and confirmed across four sign/ordering convention variants.

\subsection{Negative Audit}

A follow-up bit-structure audit tested whether the bracket structure is derivable from bit arithmetic. The audit found partial structure for Cartan generator identification but no reliable bit-derived output-index or sign prediction. The stronger claim (bit-generated su(3)) is not established.

\subsection{Claim Boundary}

\textbf{Established claims:}
\begin{itemize}
  \item Labeling rigidity: natural binary ordering uniquely preserves su(3) commutation under fixed convention.
  \item Uniqueness: verified over all $8! = 40,320$ pure permutations.
  \item Convention robustness: result holds across four tested variants.
  \item Verification method: exhaustive, deterministic, reproducible.
  \item Priority date: June 1, 2026.
\end{itemize}

\textbf{Not established:}
\begin{itemize}
  \item That bit arithmetic alone derives the bracket structure.
  \item That the correspondence holds under all Lie algebra automorphisms.
  \item That any deployed system implements su(3) gauge invariance beyond this algebraic correspondence.
  \item Physics claims beyond the formal algebraic fact.
  \item That this result has practical utility; it remains a theoretical correspondence.
\end{itemize}

\subsection{Future Work}

\begin{enumerate}
  \item \textbf{BIT\_BRACKET\_RECONSTRUCTION\_A1} (conditional): if the bit-audit had returned positive results, test whether a proposed bit rule reconstructs the full bracket table without explicit generator reference. Current negative results preclude this.

  \item \textbf{Automorphism invariance}: extend the search beyond pure permutations to Lie algebra automorphisms of su(3) and verify uniqueness in that broader space.

  \item \textbf{Higher-dimensional extensions}: investigate whether 4-bit opcode spaces preserve su(4) structure, or whether the su(3) result is specific to the 3-bit cardinality match.

  \item \textbf{Computational governance application}: develop a governance substrate that exploits the labeling rigidity result, with explicit tracking of algebraic coherence preservation.
\end{enumerate}

\section*{Data and Code Availability}

The verification script is standalone, self-contained, and publicly archived:
\begin{itemize}
  \item \textbf{GitHub}: \url{https://github.com/erynoregon-rgb/LCOS-Core} (commit 5af03e1)
  \item \textbf{Zenodo}: \url{https://doi.org/10.5281/zenodo.20499960}
\end{itemize}

The script requires only NumPy and produces the verification output shown in Section 5. It is intended to be reproducible and auditable by any reader with Python 3 and NumPy installed.

\section{Acknowledgments}

Computational verification performed with Claude Sonnet 4.6 (Anthropic). Numerical computation and matrix algebra via NumPy (Harris et al., 2020). The author acknowledges the foundational work of Gell-Mann and colleagues in establishing the su(3) structure of strong interactions.

\begin{thebibliography}{99}

\bibitem{GellMann1961}
Gell-Mann, M. (1961). The eightfold way: A theory of strong interaction symmetry.
\textit{California Institute of Technology}, Technical Report TML-20, Pasadena, CA.

\bibitem{Weinberg1973}
Weinberg, S. (1973). Non-abelian gauge theories of the strong interactions.
\textit{Physical Review Letters}, 31(7), 494--497.

\bibitem{Fritzsch1973}
Fritzsch, H., Gell-Mann, M., \& Leutwyler, H. (1973). Advantages of the color octet gluon picture.
\textit{Physics Letters B}, 47(5), 365--368.

\bibitem{Harris2020}
Harris, C. R., et al. (2020). Array programming with NumPy.
\textit{Nature}, 585, 357--362.

\end{thebibliography}

\appendix

\section{Standalone Verification Script}

The complete verification is performed by the standalone Python script available at the URLs listed in Section~\ref*{Data and Code Availability}. The script is named \texttt{su3\_verification.py} and requires only NumPy.

To run the script:
\begin{lstlisting}[language=bash]
python3 su3_verification.py
\end{lstlisting}

The script performs the following steps in order:
\begin{enumerate}
  \item Defines the eight Gell-Mann matrices in standard normalization.
  \item Derives structure constants from the trace formula $f_{abc} = \text{Tr}([\lambda_a, \lambda_b] \lambda_c) / (4i)$.
  \item Cross-checks derived constants against a hardcoded hand table.
  \item Tests the natural binary opcode ordering and reports the score.
  \item Enumerates all $8! = 40,320$ permutations, scores each, and reports the best and second-best.
  \item Tests the four convention variants (Section~\ref{sec:convention}).
  \item Tests the negative control orderings (Section~\ref{sec:negcontrols}).
  \item Reports the bit-structure audit results (Section~\ref{sec:bitaudit}).
  \item Outputs a complete record of the verification run.
\end{enumerate}

The script is deterministic and produces identical output across runs. All floating-point operations are performed in standard 64-bit IEEE754 precision; error bars reflect machine epsilon ($\approx 2.22 \times 10^{-16}$).

\end{document}
